\documentclass[aspectratio=169]{beamer}
\usepackage{ctex, hyperref}
\usepackage[T1]{fontenc}

% other packages
\usepackage{latexsym,amsmath,xcolor,multicol,booktabs,calligra}
\usepackage{graphicx,pstricks,listings,stackengine}

% Example metadata: replace these fields with your own report information.
\author{黄小明}
\title{PKU Beamer Theme}
\subtitle{毕业设计开题报告}
\institute{北京大学xxx学院}
\date{2020年11月24日}
\usepackage{PekingU}

% defs
\def\cmd#1{\texttt{\color{red}\footnotesize $\backslash$#1}}
\def\env#1{\texttt{\color{blue}\footnotesize #1}}
\definecolor{deepblue}{rgb}{0,0,0.5}
\definecolor{deepred}{rgb}{0.6,0,0}
\definecolor{deepgreen}{rgb}{0,0.5,0}
\definecolor{halfgray}{gray}{0.55}

\lstset{
    basicstyle=\ttfamily\small,
    keywordstyle=\bfseries\color{deepblue},
    emphstyle=\ttfamily\color{deepred},    % Custom highlighting style
    stringstyle=\color{deepgreen},
    numbers=left,
    numberstyle=\small\color{halfgray},
    xleftmargin=1.8em,
    numbersep=0.6em,
    rulesepcolor=\color{red!20!green!20!blue!20},
    frame=shadowbox,
}


\begin{document}

% Beamer's sans-serif body uses the portable CJK font selected by ctex.
\begin{frame}
    \titlepage
    \begin{figure}[htpb]
        \begin{center}
            \includegraphics[width=0.15\linewidth]{pic/PKU_logo.png}
        \end{center}
    \end{figure}
\end{frame}

\begin{frame}
    \tableofcontents[sectionstyle=show,subsectionstyle=show/shaded/hide,subsubsectionstyle=show/shaded/hide]
\end{frame}


\section{课题背景}

\begin{frame}{用Beamer很高大上？}
    \begin{itemize}[<+-| alert@+>] % 当然，除了alert，手动在里面插 \pause 也行
        \item 大家都会\LaTeX{}，好多学校都有自己的Beamer主题
        \item 中文支持请选择 Xe\LaTeX{} 编译选项
        \item Overleaf项目地址位于 \url{https://www.overleaf.com/latex/templates/pku-beamer-theme/zdnpwpfnkzcc}，可以直接使用
        \item GitHub项目地址位于 \url{https://github.com/inFaaa/PKU-Beamer-Theme}，如果有bug或者feature request可以去里面提issue
    \end{itemize}
\end{frame}

\section{研究现状}

\subsection{Beamer主题分类}

\begin{frame}
    \begin{itemize}
        \item 有一些 \LaTeX{} 自带的
        \item 有一些Tsinghua的
        \item 本模板来源自 \newline \url{https://www.latexstudio.net/archives/4051.html}
        \item 参考项目：\href{https://github.com/tuna/THU-Beamer-Theme}{\color{purple}{THU Beamer Theme}} \cite{origin}
        \item 这是n+e大佬在16-17年做的一些ppt：\href{https://github.com/Trinkle23897/oi_slides}{\color{purple}{戳我}}
    \end{itemize}
\end{frame}


\section{研究内容}

\subsection{美化主题}

\begin{frame}{主题与版式}
    \begin{itemize}
        \item 顶栏仅显示章节名称，页脚合并为一行
        \item 中文正文采用黑体，保留北大红作为主题色
        \item 默认使用 16:9 宽屏，仅在大章节切换时显示过渡页
        \item 更多该模板的功能可以参考 \url{https://www.latexstudio.net/archives/4051.html}
        \item 下面列举出了一些Beamer的用法，部分节选自 \url{https://tuna.moe/event/2018/latex/}
    \end{itemize}
\end{frame}

\subsection{如何更好地做Beamer}

\begin{frame}{Why Beamer}
    \begin{itemize}
        \item \LaTeX 广泛用于学术界，期刊会议论文模板
    \end{itemize}
    \begin{table}[h]
        \centering
        \small
        \renewcommand{\arraystretch}{1.1}
        \begin{tabular*}{0.94\linewidth}{@{\extracolsep{\fill}}cc@{}}
            \toprule
            Microsoft\textsuperscript{\textregistered}  Word & \LaTeX \\
            \midrule
            文字处理工具 & 专业排版软件 \\
            容易上手，简单直观 & 容易上手 \\
            所见即所得 & 所见即所想，所想即所得 \\
            高级功能不易掌握 & 进阶难，但一般用不到 \\
            处理长文档需要丰富经验 & 和短文档处理基本无异 \\
            花费大量时间调格式 & 无需担心格式，专心作者内容 \\
            公式排版差强人意 & 尤其擅长公式排版 \\
            二进制格式，兼容性差 & 文本文件，易读、稳定 \\
            付费商业许可 & 自由免费使用 \\
            \bottomrule
        \end{tabular*}
    \end{table}
\end{frame}

\begin{frame}{排版举例}
    \begin{exampleblock}{无编号公式} % 加 * 
        \begin{equation*}
            J(\theta) = \mathbb{E}_{\pi_\theta}[G_t] = \sum_{s\in\mathcal{S}} d^\pi (s)V^\pi(s)=\sum_{s\in\mathcal{S}} d^\pi(s)\sum_{a\in\mathcal{A}}\pi_\theta(a|s)Q^\pi(s,a)
        \end{equation*}
    \end{exampleblock}
    \begin{exampleblock}{多行多列公式\footnote{如果公式中有文字出现，请用 $\backslash$mathrm\{\} 或者 $\backslash$text\{\} 包含，不然就会变成 $clip$，在公式里看起来比 $\mathrm{clip}$ 丑非常多。}}
        % 使用 & 分隔
        \begin{align}
            Q_\mathrm{target}&=r+\gamma Q^\pi(s^\prime, \pi_\theta(s^\prime)+\epsilon)\\
            \epsilon&\sim\mathrm{clip}(\mathcal{N}(0, \sigma), -c, c)\nonumber
        \end{align}
    \end{exampleblock}
\end{frame}

\begin{frame}
    \begin{exampleblock}{编号多行公式}
        % Taken from Mathmode.tex
        \begin{multline}
            A=\lim_{n\rightarrow\infty}\Delta x\left(a^{2}+\left(a^{2}+2a\Delta x+\left(\Delta x\right)^{2}\right)\right.\label{eq:reset}\\
            +\left(a^{2}+2\cdot2a\Delta x+2^{2}\left(\Delta x\right)^{2}\right)\\
            +\left(a^{2}+2\cdot3a\Delta x+3^{2}\left(\Delta x\right)^{2}\right)\\
            +\ldots\\
            \left.+\left(a^{2}+2\cdot(n-1)a\Delta x+(n-1)^{2}\left(\Delta x\right)^{2}\right)\right)\\
            =\frac{1}{3}\left(b^{3}-a^{3}\right)
        \end{multline}
    \end{exampleblock}
\end{frame}

\begin{frame}{图形与分栏}
    % From thuthesis user guide.
    \begin{columns}[c,onlytextwidth]
    \begin{column}{0.44\textwidth}
        \centering
        \psset{unit=0.72cm}
        \begin{pspicture}(-1.75,-3)(3.25,4)
            \psline[linewidth=0.25pt](0,0)(0,4)
            \rput[tl]{0}(0.2,2){$\vec e_z$}
            \rput[tr]{0}(-0.9,1.4){$\vec e$}
            \rput[tl]{0}(2.8,-1.1){$\vec C_{ptm{ext}}$}
            \rput[br]{0}(-0.3,2.1){$\theta$}
            \rput{25}(0,0){%
            \psframe[fillstyle=solid,fillcolor=lightgray,linewidth=.8pt](-0.1,-3.2)(0.1,0)}
            \rput{25}(0,0){%
            \psellipse[fillstyle=solid,fillcolor=yellow,linewidth=3pt](0,0)(1.5,0.5)}
            \rput{25}(0,0){%
            \psframe[fillstyle=solid,fillcolor=lightgray,linewidth=.8pt](-0.1,0)(0.1,3.2)}
            \rput{25}(0,0){\psline[linecolor=red,linewidth=1.5pt]{->}(0,0)(0.,2)}
%           \psRotation{0}(0,3.5){$\dot\phi$}
%           \psRotation{25}(-1.2,2.6){$\dot\psi$}
            \psline[linecolor=red,linewidth=1.25pt]{->}(0,0)(0,2)
            \psline[linecolor=red,linewidth=1.25pt]{->}(0,0)(3,-1)
            \psline[linecolor=red,linewidth=1.25pt]{->}(0,0)(2.85,-0.95)
            \psarc{->}{2.1}{90}{112.5}
            \rput[bl](.1,.01){C}
        \end{pspicture}
    \end{column}
    \begin{column}{0.5\textwidth}
        \centering
        \includegraphics[width=\linewidth,height=.66\textheight,keepaspectratio]{pic/dtmf.pdf}
    \end{column}
    \end{columns}
\end{frame}

\begin{frame}[fragile]{\LaTeX{} 常用命令}
    \begin{columns}[T,onlytextwidth]
    \begin{column}{0.31\textwidth}
    \begin{exampleblock}{章节与排版}
        \centering
        \footnotesize
        \begin{tabular}{@{}ll@{}}
            \cmd{chapter} & 章 \\
            \cmd{section} & 节 \\
            \cmd{subsection} & 小节 \\
            \cmd{paragraph} & 带题头段落 \\\hline
            \cmd{centering} & 居中对齐 \\
            \cmd{emph} & 强调 \\
            \cmd{verb} & 原样输出 \\
            \cmd{url} & 超链接 \\
        \end{tabular}
    \end{exampleblock}
    \end{column}
    \begin{column}{0.35\textwidth}
    \begin{exampleblock}{图表与引用}
        \centering
        \footnotesize
        \begin{tabular}{@{}p{.60\linewidth}@{\hspace{.03\linewidth}}p{.37\linewidth}@{}}
            \cmd{footnote} & 脚注 \\
            \cmd{item} & 列表条目 \\
            \cmd{caption} & 标题 \\
            \cmd{includegraphics} & 插入图片 \\\hline
            \cmd{label} & 标号 \\
            \cmd{cite} & 引用文献 \\
            \cmd{ref} & 交叉引用 \\
        \end{tabular}
    \end{exampleblock}
    \end{column}
    \begin{column}{0.28\textwidth}
    \begin{exampleblock}{环境}
        \centering
        \footnotesize
        \begin{tabular}{@{}ll@{}}
            \env{table} & 表格 \\
            \env{figure} & 图片 \\
            \env{equation} & 公式 \\\hline
            \env{itemize} & 无编号列表 \\
            \env{enumerate} & 编号列表 \\
            \env{description} & 描述 \\\hline
        \end{tabular}
    \end{exampleblock}
    \end{column}
    \end{columns}
\end{frame}

\begin{frame}[fragile]{\LaTeX{} 环境命令举例}
    \begin{columns}[T,onlytextwidth]
    \begin{column}{0.47\textwidth}
\begin{lstlisting}[language=TeX,basicstyle=\ttfamily\footnotesize]
\begin{itemize}
  \item A \item B
  \item C
  \begin{itemize}
    \item C-1
  \end{itemize}
\end{itemize}
\end{lstlisting}
        \medskip
        \begin{itemize}
            \item A
            \item B
            \item C
            \begin{itemize}
                \item C-1
            \end{itemize}
        \end{itemize}
    \end{column}
    \pause
    \begin{column}{0.47\textwidth}
\begin{lstlisting}[language=TeX,basicstyle=\ttfamily\footnotesize]
\begin{enumerate}
  \item 巨佬 \item 大佬
  \item 萌新
  \begin{itemize}
    \item[Infaaa] 瑟瑟发抖
  \end{itemize}
\end{enumerate}
\end{lstlisting}
        \medskip
        \begin{enumerate}
            \item 巨佬
            \item 大佬
            \item 萌新
            \begin{itemize}
                \item[Infaaa] 瑟瑟发抖
            \end{itemize}
        \end{enumerate}
    \end{column}
    \end{columns}
\end{frame}

\begin{frame}[fragile]{\LaTeX{} 数学公式}
    \begin{columns}[c,onlytextwidth]
        \begin{column}{.55\textwidth}
\begin{lstlisting}[language=TeX]
$V = \frac{4}{3}\pi r^3$

\[
  V = \frac{4}{3}\pi r^3
\]

\begin{equation}
  \label{eq:vsphere}
  V = \frac{4}{3}\pi r^3
\end{equation}
\end{lstlisting}
        \end{column}
        \begin{column}{.4\textwidth}
            $V = \frac{4}{3}\pi r^3$
            \[
                V = \frac{4}{3}\pi r^3
            \]
            \begin{equation}
                \label{eq:vsphere}
                V = \frac{4}{3}\pi r^3
            \end{equation}
        \end{column}
    \end{columns}
    \begin{itemize}
        \item 更多内容请看 \href{https://zh.wikipedia.org/wiki/Help:数学公式}{\color{purple}{这里}}
    \end{itemize}
\end{frame}

\begin{frame}[fragile]{\LaTeX{} 表格与交叉引用}
    \begin{columns}[c,onlytextwidth]
        \column{.56\textwidth}
\begin{lstlisting}[language=TeX,basicstyle=\ttfamily\footnotesize]
\begin{table}[htbp]
  \caption{编号与含义}
  \label{tab:number}
  \centering
  \begin{tabular}{cl}
    \toprule
    编号 & 含义 \\
    \midrule
    1 & 4.0 \\
    2 & 3.7 \\
    \bottomrule
  \end{tabular}
\end{table}
公式~(\ref{eq:vsphere}) 的
编号与含义请参见
表~\ref{tab:number}。
\end{lstlisting}
        \column{.4\textwidth}
        \begin{table}[htpb]
            \centering
            \caption{编号与含义}
            \label{tab:number}
            \begin{tabular}{cl}\toprule
                编号 & 含义 \\\midrule
                1 & 4.0\\
                2 & 3.7\\\bottomrule
            \end{tabular}
        \end{table}
        \normalsize 公式~(\ref{eq:vsphere})的编号与含义请参见表~\ref{tab:number}。
    \end{columns}
\end{frame}

\begin{frame}{作图}
    \begin{columns}[c,onlytextwidth]
    \begin{column}{0.66\textwidth}
    \begin{itemize}
        \item 矢量图 eps, ps, pdf
        \begin{itemize}
            \item METAPOST, pstricks, pgf $\ldots$
            \item Xfig, Dia, Visio, Inkscape $\ldots$
            \item Matlab / Excel 等保存为 pdf
        \end{itemize}
        \item 位图 png, jpg, tiff $\ldots$
        \begin{itemize}
            \item 提高清晰度，避免发虚
            \item 应尽量避免使用
        \end{itemize}
    \end{itemize}
    \end{column}
    \begin{column}{0.28\textwidth}
    \begin{figure}[htpb]
        \centering
        \includegraphics[width=0.8\linewidth]{pic/PKU_logo.png}
        \caption{PNG 校徽示例}
    \end{figure}
    \end{column}
    \end{columns}
\end{frame}

\section{计划进度}
\begin{frame}
    \begin{itemize}
        \item 一月：完成文献调研
        \item 二月：复现并评测各种Beamer主题美观程度
        \item 三、四月：美化PKU Beamer主题
        \item 五月：论文撰写
    \end{itemize}
\end{frame}

\section{参考文献}

\begin{frame}[allowframebreaks]
    \bibliography{ref}
    \bibliographystyle{alpha}
    % 如果参考文献太多的话，可以像下面这样调整字体：
    % \tiny\bibliographystyle{alpha}
\end{frame}

\begin{frame}
    \begin{center}
        {\Huge\calligra Thanks!}
    \end{center}
\end{frame}

\end{document}
